%% Explicit first-order predicate doctrines with quantifiers and equality.

\subsection{Predicates, substitution, and quantifier adjunctions}
\label{sec:p10-doctrine-quantifiers}

For a fixed context, predicates form an ordered algebra: entailment
is its order, and conjunction is its meet.  A substitution changes
the context and pulls predicates back to it.  Quantifiers are
characterized by how entailment interacts with that pullback.
The following definition states the precise first-order convention
used here.

Lawvere's original hyperdoctrine is category-valued, with cartesian
closed categories of attributes and adjoints to substitution along
all terms.  His substitution functors are not required to preserve
implication in that generality
\cite[pp.~11--13]{Lawvere2006AdjointnessReprint}.
The posetal first-order convention below uses the indexed orders,
Heyting operations, projection adjoints, and equality axioms
specified separately in
\cite[Definitions~5.2.1, 5.4.1, 5.5.1, and 5.6.1,
pp.~41--56]{Pitts1995Draft}.
Their assumptions are part of this presentation.

\begin{hbdefinition}{def:p10-first-order-hyperdoctrine}{A first-order posetal hyperdoctrine}
Let \(\mathcal C\) be a locally small category with chosen finite
products.  Its objects and indexed families are formed in the
containing ambient universe of
Section~\ref{sec:p10-categorical-toolkit}.
Let \(\mathsf{Heyt}\) have small Heyting algebras as objects and
maps preserving \(0,1,\wedge,\vee,\Rightarrow\) as arrows.
An indexed predicate algebra is a functor
\[
                       P:\mathcal C^{\mathrm{op}}\longrightarrow
                         \mathsf{Heyt}.
\]
The Heyting algebra \(P(I)\) is the fibre of predicates in context
\(I\).  Write its greatest and least elements as
\(\top_I,\bot_I\).
For \(f:J\to I\), write
\[
                    f^*=P(f):P(I)\longrightarrow P(J).
\]
This reindexing preserves all the stated operations and satisfies
\(1_I^*=1_{P(I)}\) and \((f g)^*=g^*f^*\).

For every object \(A\) and product projection
\(\pi_{I,A}:I\times A\to I\), require monotone maps
\[
        \exists_{\pi_{I,A}},\forall_{\pi_{I,A}}:
                    P(I\times A)\longrightarrow P(I)
\]
with the order adjunctions
\[
 \begin{aligned}
 \exists_\pi\alpha\leq\beta
       &\quad\Longleftrightarrow\quad\alpha\leq\pi^*\beta,\\
 \pi^*\beta\leq\alpha
       &\quad\Longleftrightarrow\quad\beta\leq\forall_\pi\alpha.
 \end{aligned}
\]
Here \(\alpha\in P(I\times A)\), \(\beta\in P(I)\).
For \(f:J\to I\), require the two Beck--Chevalley identities
\[
 \begin{aligned}
 f^*\exists_{\pi_{I,A}}
     &=\exists_{\pi_{J,A}}(f\times1_A)^*,\\
 f^*\forall_{\pi_{I,A}}
     &=\forall_{\pi_{J,A}}(f\times1_A)^*.
 \end{aligned}
\]
They express stability under parameter substitution in the product
square with upper arrow \(f\times1_A:J\times A\to I\times A\)
and lower arrow \(f:J\to I\).
These squares are pullbacks using only finite products.
No adjoints along arbitrary base arrows, or other base pullbacks,
are included in this definition.
\index{hyperdoctrine!first-order posetal}
\index{predicate!indexed algebra}\index{reindexing}
\index{Beck--Chevalley condition}
\end{hbdefinition}

\begin{hbtheorem}{thm:p10-doctrine-quantifier-laws}{proved=annex:p10.quant.1}{Order laws of the two quantifiers}
The parameter squares in the preceding definition are pullbacks.
For each projection, both adjoints are uniquely determined by
reindexing.  They satisfy
\[
 \begin{gathered}
 \alpha\leq\pi^*\exists_\pi\alpha,\qquad
 \exists_\pi\pi^*\beta\leq\beta,\\
 \pi^*\forall_\pi\alpha\leq\alpha,\qquad
 \beta\leq\forall_\pi\pi^*\beta .
 \end{gathered}
\]
Existential quantification preserves finite joins and bottom,
while universal quantification preserves finite meets and top:
\[
 \begin{aligned}
 \exists_\pi(\alpha\vee\alpha')
       &=\exists_\pi\alpha\vee\exists_\pi\alpha',
 &\exists_\pi\bot_{I\times A}&=\bot_I,\\
 \forall_\pi(\alpha\wedge\alpha')
       &=\forall_\pi\alpha\wedge\forall_\pi\alpha',
 &\forall_\pi\top_{I\times A}&=\top_I.
 \end{aligned}
\]
If the fibres are Boolean algebras, then
\[
                   \forall_\pi\alpha
                         =\neg\exists_\pi\neg\alpha .
\]
The adjunction hypotheses alone supply these order laws;
the two Beck--Chevalley equations are additional substitution
conditions in the definition.
\end{hbtheorem}

\begin{hbtheorem}{thm:p10-doctrine-frobenius}{proved=annex:p10.quant.2}{Frobenius from implication-preserving reindexing}
Let \(H,K\) be Heyting algebras, and let monotone maps
\(L:H\to K\), \(R:K\to H\) satisfy \(L\dashv R\).
If \(R\) preserves Heyting implication, then
\[
                         L(a\wedge R(b))=L(a)\wedge b
                         \quad(a\in H,\ b\in K).
\]
In particular, the first-order doctrine above satisfies
Frobenius reciprocity:
\[
             \exists_\pi(\alpha\wedge\pi^*\beta)
                    =(\exists_\pi\alpha)\wedge\beta.
\]
This conclusion uses preservation of implication.
In a presentation with only finite-meet fibres, Frobenius
must instead be required or proved from other assumptions.
\index{Frobenius reciprocity}
\end{hbtheorem}

The two projection adjunctions and their parameter naturality are
equations (5.6)--(5.9), pp.~52--53, of
\cite{Pitts1995Draft}.  Frobenius is equation (5.10) and
Remark~5.5.2 on p.~53.  The source treats these assumptions
modularly; the present Heyting-preserving convention implies
Frobenius by the displayed theorem.

\subsection{Equality with parameters}
\label{sec:p10-doctrine-equality}

\begin{hbdefinition}{def:p10-doctrine-equality}{Parameterized equality predicates}
For each object \(A\), specify a predicate
\(\delta_A\in P(A\times A)\).
For any parameter object \(I\), use the left-associated product
\((I\times A)\times A\).
Write \(\pi_I:I\times A\to I\) and \(\pi_A:I\times A\to A\).
Define
\[
 \begin{aligned}
 e_{I,A}
    &=\langle1_{I\times A},\pi_A\rangle:
                        I\times A\longrightarrow(I\times A)\times A,\\
 p_{I,A}
    &=\pi_1:(I\times A)\times A\longrightarrow I\times A,\\
 q_{I,A}
    &=\langle\pi_I\pi_1,\pi_2\rangle:
                        (I\times A)\times A\longrightarrow I\times A,\\
 r_{I,A}
    &=\langle\pi_A\pi_1,\pi_2\rangle:
                        (I\times A)\times A\longrightarrow A\times A.
 \end{aligned}
\]
Thus \(e\) duplicates the last variable, \(p\) retains its first
copy, \(q\) retains its second copy, and \(r\) selects the two
copies while discarding the parameters.
Require that the monotone map
\[
 L_e:P(I\times A)\longrightarrow P((I\times A)\times A),
 \qquad L_e(\alpha)=p^*\alpha\wedge r^*\delta_A
\]
be left adjoint to \(e^*\).  Explicitly,
\[
              p^*\alpha\wedge r^*\delta_A\leq\beta
                  \quad\Longleftrightarrow\quad
              \alpha\leq e^*\beta .
\]
This is required for every parameter object \(I\).
An equality predicate is an element of a fibre; a representation
by actual subobjects is not assumed here.
\index{equality!in an indexed predicate algebra}
\index{diagonal!equality adjunction}
\end{hbdefinition}

The parameterized adjunction is equation (5.12) and
Definition~5.6.1, pp.~55--56, of
\cite{Pitts1995Draft}.
The coordinate maps above state it for the chosen left-associated
products.  Rebracketing uses the canonical product isomorphisms.

\begin{hbtheorem}{thm:p10-doctrine-equality-laws}{proved=annex:p10.quant.3}{Reflexivity and replacement of equals}
Let \(\Delta_A=\langle1_A,1_A\rangle:A\to A\times A\) and
\(\tau_A=\langle\pi_2,\pi_1\rangle:A\times A\to A\times A\).
The parameterized equality axiom implies
\[
                    \Delta_A^*\delta_A=\top_A,\qquad
                    \tau_A^*\delta_A=\delta_A.
\]
For every parameter object \(I\) and \(\gamma\in P(I\times A)\),
\[
                     p_{I,A}^*\gamma\wedge r_{I,A}^*\delta_A
                                  \leq q_{I,A}^*\gamma .
\]
In particular, in the fibre over \((A\times A)\times A\),
\[
                     p_{A,A}^*\delta_A\wedge r_{A,A}^*\delta_A
                                  \leq q_{A,A}^*\delta_A,
\]
which is transitivity.  For every arrow \(f:A\to B\),
\[
                             \delta_A\leq(f\times f)^*\delta_B.
\]
Thus reindexing equality along a pair of terms supports reflexivity,
symmetry, transitivity, and substitution in predicates and terms.
The assertion is internal equality preservation; no reflection
of base-arrow equality is included.
\end{hbtheorem}

\subsection{Inhabited sorts and unused variables}
\label{sec:p10-doctrine-inhabited}

\begin{hbdefinition}{def:p10-inhabited-sort}{An internally inhabited sort}
An object \(A\) is inhabited in this doctrine when
\[
                  \exists_{\pi_{1,A}}\top_{1\times A}=\top_1,
\]
where \(1\) is the chosen terminal object of \(\mathcal C\).
This is an assertion about predicates in the terminal context.
A chosen global element \(a:1\to A\) is a different datum.
\index{inhabited sort}\index{nonempty domain!categorical semantics}
\end{hbdefinition}

\begin{hbtheorem}{thm:p10-doctrine-inhabited}{proved=annex:p10.quant.4}{Inhabitance and reflection along context extension}
An object \(A\) is inhabited exactly when
\(\exists_{\pi_{I,A}}\top_{I\times A}=\top_I\) for every
parameter object \(I\).
A global element \(1\to A\) suffices for inhabitance.

For an inhabited \(A\), every projection \(\pi=\pi_{I,A}\)
satisfies
\[
                   \exists_\pi\pi^*\beta=\beta
                   =\forall_\pi\pi^*\beta
                         \quad(\beta\in P(I)),
\]
and reindexing along \(\pi\) reflects order:
\[
                    \pi^*\alpha\leq\pi^*\beta
                         \quad\Longleftrightarrow\quad
                    \alpha\leq\beta.
\]
Consequently, an unused variable of an inhabited sort may be
removed from the retained variable context of a semantic sequent.
For several such variables the argument applies one extension
at a time.  Without inhabitance this removal is not supplied by
the doctrine axioms.
\end{hbtheorem}

The powerset construction below permits empty sorts and makes
this distinction concrete.  The subsequent soundness theorem
retains variable contexts, then states the additional bridge to
ordinary nonempty-domain first-order semantics.
